\documentclass[11pt,a4paper,openright,twoside]{book}

% F-19: common_preamble already loads morewrites, inputenc, fontenc, the
% amsmath family, mathrsfs, microtype, booktabs, tabularx, array, xcolor,
% hyperref, geometry, enumitem, graphicx, cite, silence, tikz, plus the
% theorem environments and auto-harvested macros every chapter body needs
% (see common_preamble.tex). Only the two packages below are NOT in it:
% pdfpages (for the Dictionary's \includepdf) and emptypage (blank-page
% headers), both specific to this top-level master.
% F-19 prototype: single shared preamble for all chapter bodies.
% Union of packages/macros/theorem environments found across chap01-13
% (see INVENTARIO.md). Macros use \providecommand so the FIRST definition
% wins; the handful of cosmetic variants found by inventory.py
% (\mathrm{BV} vs \operatorname{BV}, etc.) render identically.

% morewrites MUST load before hyperref/other packages claim \write streams:
% 13 chapters x (bibliography + hyperref + float aux) exhaust the classic
% 16-stream limit; morewrites raises the ceiling to 9000+ via etex registers.
\usepackage{morewrites}
\usepackage[utf8]{inputenc}
\usepackage[T1]{fontenc}
\usepackage{amsmath,amssymb,amsfonts,amsthm,mathtools}
\usepackage{mathrsfs}
\usepackage{microtype}
\usepackage{booktabs}
\usepackage{tabularx}
\usepackage{array}
\usepackage[dvipsnames]{xcolor}
\usepackage{hyperref}
\usepackage{geometry}
\usepackage{enumitem}
\usepackage{graphicx}
% chap13's \includegraphics uses a bare filename (figure lives next to the
% chapter source, one level up from this prototype folder); point graphicx
% there instead of editing the chapter body.
\graphicspath{{../}}
\usepackage{cite}
\usepackage{silence}
\WarningFilter{nameref}{The definition of \label has changed!}
\usepackage{tikz}
% chap13 is revtex4-2; extract.py rewrites its \ruledtabular table into a
% plain booktabs \toprule/\bottomrule + center (see extract.py), so no
% \ruledtabular shim is needed here.
% revtex4-2's \acknowledgments environment (chap13) is just a run-in
% "Acknowledgments." heading; book has no built-in equivalent.
\newenvironment{acknowledgments}
  {\par\medskip\noindent\textbf{Acknowledgments.}\ }
  {\par}

\numberwithin{equation}{chapter}

% Paragraph reflow: each chapter was authored against its OWN margin
% (0.9in for ch. 3-6, 1in for ch. 7-9, a custom \geometry for ch. 1-2,
% class defaults for ch. 10-13) before being unified here to 3cm/2.8cm
% (the real master's own geometry). Raising tolerance/emergency stretch
% is the standard, honest way to let TeX re-justify those paragraphs at
% the new width instead of flagging every minor over/underfull line;
% it changes glue stretch, not the visible word choice.
\tolerance=1000
\emergencystretch=2em

% F-19 adoption: attempt to clear the remaining overfull \hbox WITHOUT
% editing any chapter. \allowdisplaybreaks lets long multi-line align/
% gather blocks break across pages instead of being forced into a cramped
% single page. A first attempt also tried \everydisplay{\small} (a
% smaller display-math font); REVERTED -- \small is a text-mode size
% switch and \everydisplay inserts it already inside math mode, which
% raised a "Command \small invalid in math mode" LaTeX Font Warning on
% every single display equation in the book (dozens), failing the
% 0-warnings gate outright, and even made the worst remaining overfull
% box bigger (64pt vs. 26pt) through inconsistent math sizing. Any
% working display-font-shrink would need a different hook (e.g.
% \AtBeginEnvironment{equation}{\small} from etoolbox, repeated for every
% display environment the chapters use) and is left for a future pass;
% see ADOCAO_F19.md for the resulting overfull count with just
% \allowdisplaybreaks plus the existing \tolerance/\emergencystretch.
\allowdisplaybreaks[4]

% --- Theorem environments (union across chapters; numbered per chapter) ---
\theoremstyle{plain}
\newtheorem{theorem}{Theorem}[chapter]
\newtheorem{lemma}[theorem]{Lemma}
\newtheorem{proposition}[theorem]{Proposition}
\newtheorem{corollary}[theorem]{Corollary}
\newtheorem{conjecture}[theorem]{Conjecture}
\newtheorem{problem}[theorem]{Problem}
\newtheorem{algorithm}[theorem]{Algorithm}

\theoremstyle{definition}
\newtheorem{definition}[theorem]{Definition}
\newtheorem{example}[theorem]{Example}
\newtheorem{construction}[theorem]{Construction}
\newtheorem{assumption}[theorem]{Assumption}

\theoremstyle{remark}
\newtheorem{remark}[theorem]{Remark}

% --- Custom notation: AUTO-HARVESTED from every chapter's preamble by
% extract.py (union, first definition wins, each as \providecommand).
% See gen/common_macros.tex. Do not hand-list macros here: re-run
% extract.py after touching a chapter's \newcommand/\providecommand block.
\input{gen/common_macros}

% 2026-10-11: \xref{label}{fallback} is harvested above as a standalone
% fallback (\providecommand{\xref}[2]{#2}, defined in whichever chapter's
% preamble extract.py scans first -- see chap05/09/11/12). In the master,
% EVERY chapter shares one numbering, so override it here to resolve the
% real cross-chapter number via \ref; #1 is already the target label with
% its chapNN: prefix (e.g. chap07:thm:lower_bounds), untouched by
% extract.py's own \ref-prefixing because \xref is not in its REF_CMDS list.
\renewcommand{\xref}[2]{\ref{#1}}

% book's default "headings" pagestyle puts the CURRENT SECTION title
% (via \sectionmark/\rightmark) in the running head on right-hand pages.
% Several chapters have section titles far longer than the text width
% (e.g. "2.6. TENSOR NETWORKS: GAUGE THEORY, PATH INTEGRALS, AND CONTINUOUS
% LIMITS"), which overflows the head rule regardless of \chapter[short]{}.
% Drop section titles from the running head; keep the chapter's short
% title (set via \chaptermark in the generated master) on both sides.
\renewcommand{\sectionmark}[1]{}

% --- Manual "Abstract" block (book class has no native abstract env;
% \quotation/\endquotation is a standard LaTeX kernel environment) ---
\newenvironment{chapterabstract}
  {\begin{quotation}\noindent\textbf{Abstract.}\ }
  {\end{quotation}}

\usepackage{pdfpages}
\usepackage{emptypage}
\usepackage{longtable}

\geometry{
  top=3cm,
  bottom=3cm,
  left=3cm,
  right=2.8cm,
  headheight=14pt
}

\hypersetup{
  colorlinks=true,
  linkcolor=NavyBlue,
  citecolor=Mahogany,
  urlcolor=TealBlue,
  pdftitle={Geometry, Tensors, and Quantum Gravity: Foundations of Simplicial Geometry, Minimax Foliations, and Emergent Spacetime},
  pdfauthor={Reinaldo M. Silva-Filho}
}

\begin{document}

\frontmatter
\begin{titlepage}
    \centering
    \vspace*{2cm}
    
    {\Huge \textbf{Geometry, Tensors, and Quantum Gravity} \par}
    \vspace{0.6cm}
    {\LARGE \textbf{Foundations of Simplicial Geometry, Minimax Foliations,} \\ \textbf{and Emergent Spacetime} \par}
    \vspace{1cm}
    {\large \textit{From Simplicial Fractional Calculus and Minimax Foliations} \\ \textit{to Emergent Holographic Spacetime and Observational Signatures} \par}
    
    \vspace{3cm}
    
    {\Large \textbf{Reinaldo M. Silva-Filho} \par}
    \vspace{0.5cm}
    {\normalsize Programa de P\'os-Gradua\c{c}\~ao em Estat\'istica e Experimenta\c{c}\~ao Agropecu\'aria (PPGEE/DES) \par}
    {\normalsize Departamento de Estat\'istica (DES) \par}
    {\normalsize Universidade Federal de Lavras (UFLA) \par}
    
    \vspace{0.8cm}
    {\normalsize ORCID: \href{https://orcid.org/0009-0003-8068-3330}{0009-0003-8068-3330} \par}
    {\normalsize E-mail: \href{mailto:reinaldo.filho1@estudante.ufla.br}{reinaldo.filho1@estudante.ufla.br} \par}
    
    \vspace{1.5cm}
    {\normalsize \textit{Supported by Coordena\c{c}\~ao de Aperfei\c{c}oamento de Pessoal de N\'ivel Superior} \par}
    {\normalsize \textit{(CAPES) -- Finance Code 001} \par}
    
    \vfill
    {\large \today \par}
\end{titlepage}

\vspace*{0.2\textheight}
\begin{flushright}
\textit{"Either mathematics is too big for the human mind, \\
or the human mind is more than a machine."} \\
\vspace{0.2cm}
--- Kurt G\"odel
\end{flushright}

\vspace{1.5cm}

\begin{flushright}
\textit{"Mathematics resides in the dialectic between the discrete and the continuous."} \\
\vspace{0.2cm}
--- Reinaldo M. Silva-Filho
\end{flushright}
\vfill
\clearpage

\chapter*{Author's Preface}
This monograph collects work on three programs: functional tensor algebra and geometric flows, fractional calculus on continuous Pascal simplices, and the variational theory of minimax extrinsic curvature with causal Jordan embeddings.

The objective of this work is to develop the mathematics of these three programs rigorously, and to state precisely what they do and do not imply for the problem of reconciling General Relativity with Quantum Field Theory. Throughout, statements labelled \emph{Theorem}, \emph{Proposition} or \emph{Lemma} are proved in the text or cited as classical results. Statements whose proof is not available are labelled \emph{Conjecture}. Phenomenological parametrizations and heuristic analogies are identified as such, and Chapter~13 records order-of-magnitude estimates of the observational signatures discussed in Part~IV.

\section*{Principal Results}
Most of the classical material is cited. The results proved in the text that belong to this book include: the Fourier symbol of the simplicial Beta-kernel Laplacian and its long-wavelength moment matrix (Chapter~3); the continuous Star of David identity, the golden-ratio asymptotics of the continuous Fibonacci diagonal, and a tail estimate for the continuous Dixon integral at integer order (Chapter~3); the exact Sobolev gain $\gamma = 1$ of the Beta-kernel trace, an inversion formula from the family of traces, and the sharp Lipschitz constant $\alpha/(\alpha+m)$ of the induced barycentric map (Chapter~5); the knot gap $4\pi/L$, the contact principle at obstacles, the structure of the minimiser of the linear wall problem, and the bound $\kappa R \ge 1$ for closed submanifolds in a ball, which rules out corrugated competitors of Nash--Kuiper type (Chapter~7); U-turn bounds in tubes of the space forms and pointwise bounds on slices of bounded extrinsic curvature (Chapters~8 and~12); the minimality of the three-arc teardrop U-turn among three-arc curves through a narrow mouth (Chapter~9); and a counterexample to curvature-based generalization bounds for interpolating linear regression (Chapter~10). A formalization in \textsc{Lean}~4 against \textsc{Mathlib} covers selected results of Chapters~1--12 and the numerical bounds of Chapter~13, each in a declared form (full, finite-dimensional, pointwise or conditional on named classical theorems). The Results Index at the end of the book lists, for every theorem, proposition, lemma and corollary, its proof, its numerical checks, its \textsc{Lean} module with the declared reduction, and its audit status.

\section*{Status of the Verification Material}
The accompanying repository contains two kinds of supporting material. Their status is stated here explicitly.
{\sloppy
\begin{enumerate}
    \item \textbf{Numerical scripts.} The scripts \texttt{verify\_chapXX\_numerical.py} (one per chapter) and the auxiliary checking scripts of the repository evaluate formulas of the text numerically, for example by comparing closed forms with direct quadrature. Such checks catch algebraic and numerical errors. They do not replace proofs, and a check can only be as independent as its reference computation. The checks compare against an independent computation and include a negative control, a mutated formula that must fail.
    \item \textbf{\textsc{Lean 4} files} (\texttt{formal\_proofs\_book}, \texttt{formal\_proofs\_lean4}). The modules under \texttt{formal\_proofs\_book/Book} are a \emph{naming skeleton}: each chapter's results are recorded as named structures whose fields are Boolean or floating-point placeholders, and the accompanying Lean ``theorems'' only restate their hypotheses. They therefore verify \emph{no} mathematical content of this book and should not be read as a formalization. A separate formalization against \textsc{Mathlib} covers selected results, listed with their axiom check in \texttt{formal\_proofs\_mathlib/README.md}; the chapters concerned and the level of each formalization are stated under ``Verification status'' in the Declarations below. Only results stated with their \textsc{Lean} statement and axiom check in a published repository should be regarded as machine-checked.
\end{enumerate}
\par}

{\sloppy
The source files, numerical scripts, and \textsc{Lean} material are hosted at
\begin{center}
\url{https://github.com/reinaldomsilvafilho-netizen/quantum-gravity-lean4}
\end{center}
together with the verification records of each result.
}

% Declarations: single source. The text lives in ../declarations_chapter.tex
% (shared with the standalone chapter PDFs); extract.py drops the per-chapter
% copies from gen/chapNN_body.tex, so the book has exactly one Declarations section.
\newcommand{\gtqgdeclscope}{book}
% Shared declarations block: the single source for the declarations of the book.
% \input by every standalone chapter before its bibliography, and \input ONCE in the
% front matter of the single-source master (official_master_structure.tex); the
% chapter bodies of the master do not repeat it (extract.py drops their \input).
% \gtqgdeclscope is "chapter" by default; the master sets it to "book" before \input.
% Aligned with the Leiden Declaration on Artificial Intelligence and Mathematics
% (arXiv:2609.35334); see _staging/DECLARACOES_IA_PROPOSTA.md Sec.4.2 for the rationale.
\providecommand{\gtqgdeclscope}{chapter}
\section*{Declarations}

\noindent\textbf{Funding.} This research was financed in part by the Coordena\c{c}\~ao de Aperfei\c{c}oamento de Pessoal de N\'ivel Superior -- Brasil (CAPES) -- Finance Code 001.

\medskip\noindent\textbf{Competing interests.} The author declares no competing interests.

\medskip\noindent\textbf{Tool and computational resource disclosure.} In line with the Leiden Declaration on Artificial Intelligence and Mathematics (arXiv:2609.35334), the author discloses that large language model assistants (Anthropic Claude; Google Gemini, used through the Antigravity agent environment) were used extensively in preparing this work: drafting and editing text and \LaTeX{} source; proposing, expanding and revising derivations and proofs under the author's direction; writing the Python verification scripts and the \textsc{Lean 4} files; searching and checking the literature and references; and adversarial auditing of proofs, numerical claims and code. These audits found errors, some of them originally introduced with AI assistance; the errors and their corrections are recorded in \texttt{WORKPLAN.md}. Where a result builds on a specific prior idea, the author has sought to identify and cite its source, and states explicitly where a satisfactory attribution could not be established.

\medskip\noindent\textbf{Authorship and responsibility.} AI tools are not authors of this work. Credit and responsibility for its content belong solely to the human author, who chose the research questions and the overall framework, reviewed the material, and is responsible for the content, including any remaining errors.

\medskip\noindent\textbf{Verification status.} Results labelled Theorem, Proposition or Lemma are proved in the text or cited as classical; statements labelled Conjecture are not proved. The numerical scripts compare formulas of the text with independent computations and are checks, not proofs. The \textsc{Lean 4} modules in \texttt{formal\_proofs\_book/Book} are a naming skeleton and do not verify the mathematics of this \gtqgdeclscope; they do not constitute the formal verification the Leiden Declaration asks organisations to consider for automated results. Selected results of Chapters~1--13 are formalized in \textsc{Lean 4} with Mathlib in the repository folder \texttt{formal\_proofs\_mathlib}, at the level stated there (in several cases a declared special case); in particular, the Chapter~12 constraint bounds are formalized in a matrix-level form with the ADM constraint taken as a hypothesis, the Chapter~8 curvature identities only in their pointwise algebraic part, with the Gauss, Raychaudhuri and Einstein equations they use taken as hypotheses, and the Chapter~8 U-turn and causal statements in the reduced forms stated there, the Chapter~11 minimum-cut entanglement-entropy bound only in its linear-algebra half, with the tensor-network factorization taken as a hypothesis, together with the inequality $g^{\mathrm{BKM}} \geq g^{\mathrm{SLD}}$ for full-rank $\rho_0$ in its eigenbasis, and the Chapter~13 closed-form running spectral dimension and tensor-tilt arithmetic with the bracket-nonzero hypothesis of that formula left unproved (true on the whole domain by Mills' inequality, not formalized) and the CODATA inputs taken as data. The same folder also covers selected results of the \textit{Simplicial Quantum Gravity}, Yang--Mills, fermion-mass, \textit{Beyond the Spectrum} and \textit{Functorial Bridge} companion works.

\medskip\noindent\textbf{Publication status.} This monograph is distributed as a preprint on Zenodo under the CC~BY~4.0 license. It has not yet been through peer review; Zenodo is a repository, not a peer-reviewed venue.

\medskip\noindent\textbf{Code and data availability.} Sources, numerical scripts and the audit log are available at
\begin{center}\url{https://github.com/reinaldomsilvafilho-netizen/quantum-gravity-lean4}.\end{center}
The monograph is archived on Zenodo, \href{https://doi.org/10.5281/zenodo.22290043}{doi:10.5281/zenodo.22290043}, under the CC~BY~4.0 license.


\tableofcontents

\mainmatter

% ==============================================================================
\part{Foundations of Functional Algebra, Hypertensors, and Geometric Flows}
% ==============================================================================

\chapter{Functional Realizations of Matrices and Hypertensors: Emergent Invariants and Geometric Measures}
\chaptermark{Realizations of Matrices and Hypertensors}
\input{gen/chap01_body}

\chapter{Geometric Flows, Partial Differential Equations, and Variational Dynamics on Matrix and Tensor Manifolds}
\chaptermark{Geometric Flows on Matrix and Tensor Manifolds}
\input{gen/chap02_body}

% ==============================================================================
\part{Continuous Simplicial Geometry and Fractional Calculus on Pascal Simplices}
% ==============================================================================

\chapter{Analytic Continuation of Pascal's Simplex: Continuous Multinomial Integrals, Polytope Boundary Recurrences, and Simplicial Fractional Calculus}
\chaptermark{Analytic Continuation of Pascal's Simplex}
\input{gen/chap03_body}

\chapter{Nonlinear Simplicial Waves and Anomalous Porous Transport Induced by Beta-Kernel Fractional Laplacians}
\chaptermark{Nonlinear Simplicial Waves and Porous Transport}
\input{gen/chap04_body}

\chapter{Inter-Dimensional Simplicial Transforms, Fractional Boundary Traces, and Matrix Beta-Kernels}
\chaptermark{Inter-Dimensional Simplicial Transforms}
\input{gen/chap05_body}

\chapter{Continuous Pascal Simplexes, Sierpi\'nski Gasket Laplacians, and Multifractal Singularity Spectra}
\chaptermark{Sierpi\'nski Gasket Laplacians}
\input{gen/chap06_body}

% ==============================================================================
\part{Non-Euclidean Minimax Extrinsic Curvature and Relativistic Foliations}
% ==============================================================================

\chapter{Minimax-Flat \texorpdfstring{$k$}{k}-Submanifolds in Obstacle Environments: Variational Theory, Constructive Heuristics, and Explicit Candidates}
\chaptermark{Minimax-Flat Submanifolds}
\input{gen/chap07_body}

\chapter{Minimax Extrinsic Curvature in Non-Euclidean Geometries and General Relativity: Variational Theory, Space Forms, and ADM Spacetime Slicings}
\chaptermark{Minimax Curvature and ADM Slicings}
\input{gen/chap08_body}

\chapter{Global Minimax Curvature in Multiply-Connected Domains: Homotopy Search, Universal Covering Spaces, and Non-Trivial Loop Dynamics}
\chaptermark{Global Minimax Curvature}
\input{gen/chap09_body}

\chapter{Minimax Curvature Trajectories on Statistical Manifolds: A Proposal for Deep Learning and Its Limits}
\chaptermark{Minimax Trajectories on Statistical Manifolds}
\input{gen/chap10_body}

% ==============================================================================
\part{Emergent Spacetime: Established Results, Conjectures, and Heuristics}
% ==============================================================================

\chapter[Emergent Spacetime and Quantum Geometry: Tensor Networks, Holo\-no\-mies, and Entanglement (Established Results and Conjectures)]{Emergent Spacetime and Quantum Geometry: Tensor Networks, Holonomies, and Entanglement (Established Results and Conjectures)}
\chaptermark{Emergent Spacetime and Quantum Geometry}
\input{gen/chap11_body}

\chapter{Emergent Spacetime, Holographic Entanglement, and Quantum Geometry: Proved Results, Conjectures, and Observational Limits}
\chaptermark{Holographic Entanglement and Quantum Geometry}
\input{gen/chap12_body}

% ==============================================================================
\part{Observational Signatures: Order-of-Magnitude Estimates}
% ==============================================================================

\chapter{Observational Signatures of a Running Spectral Dimension: Parametrization and Order-of-Magnitude Assessment}
\chaptermark{Observational Signatures}
\input{gen/chap13_body}

\backmatter

\chapter[Comprehensive Dictionary of Terms, Symbols, and Core Formulations]{Comprehensive Dictionary of Terms, Symbols, and Core Formulations}
\chaptermark{Dictionary of Terms and Symbols}
\includepdf[pages=1-]{DICTIONARY_TERMS_AND_SYMBOLS.pdf}

% Results Index: generated by gen_results_index.py from the chapter labels and
% formal_proofs_mathlib/README.md; the file lives one level up, like the declarations.
% Generated by _single_source_prototype/gen_results_index.py -- do not edit by hand.
\chapter[Results Index and Verification Status]{Results Index and Verification Status}
\chaptermark{Results Index}

This index lists every theorem, proposition, lemma and corollary of Chapters~1--13 with the support it has. Conjectures, definitions and remarks are not listed. The columns are:
{\raggedright
\begin{itemize}[leftmargin=1.5em,itemsep=1pt]
  \item \textbf{Proof}: \emph{text} (complete proof in the chapter), \emph{sketch} (proof outline), \emph{cited} (classical result with a precise citation, with or without a proof or sketch in the text), \emph{Ch.~3} (restates a theorem proved in Chapter~3).
  \item \textbf{Num.}: numerical checks of the repository that compare the statement, or a formula in it, with an independent computation and a negative control. V: the chapter's script \texttt{verify\_chapNN\_numerical.py}; F: the chapter's independent oracle script \texttt{F02\_chNN.py}; A: another targeted checking script. A check supports a formula; it does not prove the statement.
  \item \textbf{Lean}: the module of the \textsc{Lean}~4/\textsc{Mathlib} formalization (\texttt{formal\_proofs\_mathlib/LeanReal}) and the reduction declared for it in \texttt{formal\_proofs\_mathlib/README.md}. Every listed module compiles without \texttt{sorry} or new axioms. A module formalizes the stated reduction only (for instance a finite-dimensional, pointwise or conditional form), not the full statement unless it says \emph{full}. The \textsc{Lean} files in \texttt{formal\_proofs\_book} are a naming skeleton and are not listed.
  \item \textbf{Status}: L1--L2 means that the chapter's results were checked in a blind referee report (layer~1), revised by a different hand and re-checked (layer~2). An F-number names a later finding about the result; it is closed unless marked otherwise.
\end{itemize}
\par}
A dash means none.

{\scriptsize
\setlength{\tabcolsep}{3pt}
\renewcommand{\arraystretch}{1.15}
\begin{longtable}{@{} >{\raggedright\arraybackslash}p{0.9cm} >{\raggedright\arraybackslash}p{3.7cm} >{\raggedright\arraybackslash}p{1.25cm} >{\raggedright\arraybackslash}p{0.85cm} >{\raggedright\arraybackslash}p{5.3cm} >{\raggedright\arraybackslash}p{1.55cm} @{}}
\toprule
\textbf{No.} & \textbf{Result} & \textbf{Proof} & \textbf{Num.} & \textbf{Lean module: reduction} & \textbf{Status} \\
\midrule
\endfirsthead
\toprule
\textbf{No.} & \textbf{Result} & \textbf{Proof} & \textbf{Num.} & \textbf{Lean module: reduction} & \textbf{Status} \\
\midrule
\endhead
\midrule
\multicolumn{6}{r@{}}{\emph{continued on the next page}} \\
\endfoot
\bottomrule
\endlastfoot

\multicolumn{6}{@{}l}{\textbf{Chapter 1. Functional Realizations of Matrices and Hypertensors}} \\[1pt]
\ref{chap01:prop:spectral_critical} & Prop. Spectral and Critical Correspondence & text & V & \texttt{Chap01Rayleigh}: (i), finite-dimensional; \texttt{BeyondSpectrum1\_B3}: (ii), algebraic part in eigen-coordinates & L1--L2 \\
\ref{chap01:thm:dirichlet_blindness} & Thm. Spectral Blindness to Spatial Energy & text & V, F & \texttt{BeyondSpectrum1}: algebraic part, coefficient form (Parseval step not formalized) & L1--L2 \\
\ref{chap01:thm:tv_matrix} & Thm. Total Variation of the Matrix Step Realization & text & V, F & -- & L1--L2 \\
\ref{chap01:thm:coarea_linkage} & Thm. Geometric Coarea Linkage: Hausdorff Measure of Matrix Level Curves & text & V & -- & L1--L2 \\
\ref{chap01:thm:morse_matrix} & Thm. Morse Spectrum of Quadratic Realizations & text & V & -- & L1--L2 \\
\ref{chap01:thm:grothendieck_cutnorm} & Thm. Equivalence with the $L^\infty \to L^1$ Operator Norm & text & V, A & -- & L1--L2; F-60 \\
\ref{chap01:thm:graphon_compactness} & Thm. Compactness of the Continuum Moduli Space & cited & -- & -- & L1--L2 \\
\ref{chap01:thm:tensor_weierstrass} & Thm. Functional Well-Posedness over Compact Product Manifolds & text & -- & \texttt{Chap01Tensor}: (i) full; (ii) algebraic step, multipliers as hypothesis & L1--L2 \\
\ref{chap01:thm:hypergraphon_compactness} & Thm. Multilinear Operator Duality and Hypergraphon Compactness & text & A & \texttt{Chap01Hypergraphon}: (i), finite $k$-linear case; (ii) not formalized & L1--L2; F-60 \\
\ref{chap01:thm:attention_sobolev} & Prop. Sobolev Bound for Attention Realizations & text & A & \texttt{Chap01Attention}: coefficient form (Parseval not formalized) & L1--L2 \\
\ref{chap01:thm:kac_rice_tensors} & Thm. Complexity of the Isotropic Tensor Landscape & text & A & \texttt{Chap01KacRice}: (i), (iv), algebraic core; covariance as hypothesis & L1--L2 \\
\ref{chap01:thm:critical_scaling} & Thm. Critical Multilinear Scaling and Sub-Gaussian Concentration & text & V & \texttt{Chap01Gaussian}: (ii) for fixed unit vectors; (i) not formalized & L1--L2 \\
\addlinespace[2pt]
\multicolumn{6}{@{}l}{\textbf{Chapter 2. Geometric Flows on Matrix and Tensor Manifolds}} \\[1pt]
\ref{chap02:thm:cone_curvature} & Thm. Curvature of the Affine-Invariant Cone & text & V & -- & L1--L2 \\
\ref{chap02:prop:tangent_proj_mr} & Prop. Tangent Space Projection on $\mathcal{M}_r$ & cited + sketch & V, A & \texttt{Chap02Tangent}: matrix algebra of the projector; manifold structure not formalized & L1--L2; F-60 \\
\ref{chap02:thm:tt_manifold} & Thm. Smooth Embedded Manifold Structure of Tensor-Trains & text & F & -- & L1--L2 \\
\ref{chap02:thm:tt_tangent_proj} & Thm. Tangent Space Projection and Projected Gradient Flow & text & V & \texttt{Chap02Tangent}: (iii), abstract projected flow & L1--L2 \\
\ref{chap02:thm:toda_qr} & Thm. Isospectrality and QR Interpolation & text & V, F, A & \texttt{Chap02Toda}: algebraic and pointwise-derivative steps; ODE existence, uniqueness and QR not formalized & L1--L2; F-60 \\
\ref{chap02:thm:rank_preservation} & Thm. Rank Invariance and Energy Monotonicity & text & V & \texttt{Chap02Tangent}: (ii) and its example & L1--L2 \\
\ref{chap02:prop:graphon_dirichlet_energy} & Prop. Dirichlet Energy of the Graphon Laplacian & text & V & \texttt{Chap02Graphon}: finite graph, symmetric $W$ & L1--L2 \\
\ref{chap02:thm:graphon_heat_existence} & Thm. Global Well-Posedness and Semigroup Generation & text & V, F & \texttt{Chap02Heat}: finite grid (step graphons) & L1--L2 \\
\ref{chap02:thm:heat_smoothing} & Thm. Cut-Norm Contraction and Decay of Singular Supports & text & V & \texttt{Chap02Heat}: finite grid (step graphons) & L1--L2 \\
\ref{chap02:thm:ricci_neckpinch} & Prop. Curvature-Controlled Decay and Growth & text & V & \texttt{Chap02Flows}: scalar ODE at fixed $(x,y)$, curvature a given function & L1--L2 \\
\ref{chap02:thm:mean_curvature_coarea} & Prop. Dissipation of Level-Set Perimeter under MCF & text & -- & -- & L1--L2 \\
\ref{chap02:thm:cmps_energy_flow} & Prop. Projected Imaginary-Time Evolution on cMPS; Formal Derivation & text & A & \texttt{Chap02Flows}: finite-parameter core & L1--L2 \\
\ref{chap02:thm:continuous_path_integral} & Thm. Continuum Path-Integral Limit and Non-Abelian Gauge Invariance & text & V, F, A & \texttt{Chap02PathIntegral}: (iii) and the example in (ii); bound (i)--(ii) not formalized & L1--L2; F-63 \\
\addlinespace[2pt]
\multicolumn{6}{@{}l}{\textbf{Chapter 3. Analytic Continuation of Pascal's Simplex}} \\[1pt]
\ref{chap03:prop:stifel} & Prop. Global Stifel Recurrence & text & V, F & \texttt{Chap03Pascal}: full & L1--L2 \\
\ref{chap03:prop:digamma_pde} & Prop. Digamma Partial Differential Equation System & text & V, A & \texttt{Chap03Special}: off the poles, one variable at a time & L1--L2; F-60 \\
\ref{chap03:prop:meromorphic_zero_loci} & Prop. Global Meromorphic Extension and Zero Loci & sketch & -- & \texttt{Chap03Pascal}, \texttt{Chap03Analysis}: zeros, positivity, reflection identity; meromorphic extension not formalized & L1--L2 \\
\ref{chap03:thm:2d_row_integral} & Thm. Exact Trigonometric Representation of Row Integral & text & V, F & -- & L1--L2 \\
\ref{chap03:thm:simplex_integral} & Thm. Leading Asymptotics of the Simplex Volume & text & -- & -- & L1--L2 \\
\ref{chap03:thm:star_of_david} & Thm. Continuous Star of David Theorem & text & -- & \texttt{Chap03Pascal}: full & L1--L2 \\
\ref{chap03:thm:fibonacci} & Thm. Continuous Fibonacci Scaling and Ray Integrals & text & V, F & \texttt{Chap03Special}: algebraic core; asymptotics not formalized & L1--L2 \\
\ref{chap03:thm:lp_rows} & Thm. $L^p$ Row Norms of the Continuous Pascal Triangle & text & F & -- & L1--L2 \\
\ref{chap03:conj:dixon} & Thm. Continuous Dixon Integral, integer case & text & F & -- & L1--L2; F-54 \\
\ref{chap03:cor:dixon_higher} & Cor. Higher harmonics & text & -- & -- & L1--L2; F-54 \\
\ref{chap03:cor:dixon_real} & Cor. Real $x$ & text & A & -- & L1--L2; F-54 \\
\ref{chap03:thm:alternating_row_integrals} & Thm. Alternating Row Integrals & text & V, F & \texttt{Chap03Analysis}: full & L1--L2 \\
\ref{chap03:prop:hockey} & Prop. Continuous Longitudinal Column Integral & text & -- & -- & L1--L2 \\
\ref{chap03:thm:moments} & Thm. Simplex Moments and Covariance Structure & text & V & -- & L1--L2 \\
\ref{chap03:cor:row_log_product} & Cor. Closed Form for Continuous Row Log-Product & text & A & \texttt{Chap03Analysis}: conditional on the Alexeiewsky identity (hypothesis) & L1--L2; F-56 \\
\ref{chap03:prop:fractional_props} & Prop. Properties of the Simplicial Operator & text & -- & \texttt{Chap03Special}: (1); (2) for $m = 2$; (3) for real $\lambda$ & L1--L2 \\
\ref{chap03:thm:laplacian_symbol} & Thm. Fourier Symbol and Long-Wavelength Expansion & text & A & \texttt{Chap04Symbol}: abstract probability measure; permutation invariance as hypothesis & L1--L2 \\
\addlinespace[2pt]
\multicolumn{6}{@{}l}{\textbf{Chapter 4. Nonlinear Simplicial Waves and Porous Transport}} \\[1pt]
\ref{chap04:prop:self_adjoint_positive} & Prop. Self-Adjointness and Positivity & text & V & \texttt{Chap04SelfAdjoint}: finite abelian group in place of $\mathbb{R}^{m-1}$ & L1--L2 \\
\ref{chap04:thm:dispersion} & Thm. Dispersion Symbol and Long-Wavelength Metric & Ch.~3 & V, F & \texttt{Chap04Symbol}: abstract probability measure; permutation invariance as hypothesis & L1--L2 \\
\ref{chap04:thm:nlse_conservation} & Thm. Global Conservation of Mass and Energy & text & V, F & \texttt{Chap04Conservation}: finite lattice, Hermitian matrix in place of $\Delta^\alpha$ & L1--L2 \\
\ref{chap04:thm:modulational_instability} & Thm. Simplicial Modulational Instability & text & V, A & \texttt{Chap04Modulational}: after linearization & L1--L2; F-60 \\
\ref{chap04:thm:frequency_closed_form} & Thm. Closed-Form Solution in Frequency Domain & text & V & -- & L1--L2 \\
\ref{chap04:thm:msd} & Thm. Anisotropic Mean Squared Displacement & text & V, F & \texttt{Chap04Moments}: conditional ($E_\beta'(0)$ and $\operatorname{Hess}\sigma(0)$ as hypotheses) & L1--L2 \\
\addlinespace[2pt]
\multicolumn{6}{@{}l}{\textbf{Chapter 5. Inter-Dimensional Simplicial Transforms}} \\[1pt]
\ref{chap05:prop:fourier_rep} & Prop. Fourier Representation & text & V & -- & L1--L2 \\
\ref{chap05:thm:sobolev_trace} & Thm. Sobolev Trace Theorem with Kernel Smoothing & text & V, F & -- & L1--L2 \\
\ref{chap05:thm:dual_extension} & Thm. Dual Simplicial Extension Operator & text & V & -- & L1--L2 \\
\ref{chap05:prop:coupled_conservation} & Prop. Conservation of Total Mass and Energy Dissipation; formal computation & text & V & \texttt{Chap05Coupled}: abstract inner-product spaces; Green and Neumann identities as hypotheses & L1--L2 \\
\ref{chap05:prop:inversion} & Prop. Inversion & text & V, F, A & \texttt{Chap05Barycentric}: finite-dimensional algebra; Fourier step not formalized; $\widehat{\mathcal{K}}_\alpha \ne 0$ as hypothesis & L1--L2; F-60 \\
\ref{chap05:thm:barycentric_ratio} & Thm. Barycentric Lipschitz Bound with Sharp Constant & text & V, F & \texttt{Chap05Barycentric}: full (slightly more general than the text) & L1--L2 \\
\ref{chap05:thm:siegel_wishart_invariance} & Thm. Invariance and Spherical Harmonic Eigenfunctions & cited & V, F & -- & L1--L2 \\
\addlinespace[2pt]
\multicolumn{6}{@{}l}{\textbf{Chapter 6. Sierpi\'nski Gasket Laplacians}} \\[1pt]
\ref{chap06:assump:kigami_decimation} & Prop. Harmonic Decimation Renormalization Factor & text & V, F & \texttt{Chap06Kigami}: exact for $m = 2, 3$; reduced equations for general $m$ & L1--L2 \\
\ref{chap06:thm:sierpinski_spectral_dimension} & Thm. Spectral Dimension of the Sierpi\'nski Simplex & cited & V, F & -- & L1--L2 \\
\ref{chap06:thm:multifractal} & Prop. Asymptotics of the Continuous Row Entropy & text & F & \texttt{Chap06Multifractal}: conditional on the Stirling and Barnes series (hypotheses) & L1--L2 \\
\ref{chap06:prop:pascal_mod2} & Prop. Box-counting dimension of Pascal's triangle modulo $2$ & cited + proof & F & \texttt{Chap06PascalMod2}: combinatorial part; dimension not formalized & L1--L2 \\
\addlinespace[2pt]
\multicolumn{6}{@{}l}{\textbf{Chapter 7. Minimax-Flat Submanifolds}} \\[1pt]
\ref{chap07:prop:knot_gap} & Prop. Knotting forces a curvature gap & text & F & \texttt{Chap07Bounds}: conditional on Fenchel and F\'ary--Milnor (hypotheses) & L1--L2 \\
\ref{chap07:prop:dubins_duality} & Prop. Dubins duality & text & -- & \texttt{Chap08Chronology}: ordinal reduction; Dubins' theorem as hypothesis & L1--L2 \\
\ref{chap07:lem:wall_structure} & Lem. Structure of the minimiser of $P_W^T$ & text & -- & \texttt{Chap07Control}: (c)--(d); item (b) as hypothesis & L1--L2; F-52 \\
\ref{chap07:thm:obstacle_exclusion} & Prop. Contact principle & text & V, A & \texttt{Chap07Control}: (1)--(2) for curves; distance-function facts as hypotheses & L1--L2 \\
\ref{chap07:thm:lower_bounds} & Prop. Elementary lower bounds & text & V, F & \texttt{Chap08UTurn}: (3), differentiable tangent angle & L1--L2 \\
\ref{chap07:prop:sharp_disc_spacing} & Prop. Contact spacing at sharp discs & text & -- & -- & L1--L2 \\
\ref{chap07:thm:existence} & Prop. Existence of minimax curves & text & -- & -- & L1--L2 \\
\ref{chap07:lem:smoothing_interior} & Lem. Smoothing away from the obstacle & text & -- & -- & L1--L2 \\
\ref{chap07:prop:regularity_free} & Prop. Regularity invariance for curves without obstacles & text & V & -- & L1--L2 \\
\ref{chap07:prop:optimal_regularity} & Prop. $C^2$ minimizers of the planar problem & text & -- & \texttt{Chap07Control}: core step only & L1--L2 \\
\ref{chap07:prop:bounded_derivative_rate} & Prop. Bounded derivative of the curvature & text & -- & -- & L1--L2 \\
\ref{chap07:prop:ball_bound} & Prop. Closed submanifolds cannot be squeezed & text & A & \texttt{Chap07BallBound}: closed curves, twice differentiable & L1--L2; F-52 \\
\ref{chap07:thm:monotonicity} & Thm. Dimensional Monotonicity & text & -- & \texttt{Chap07Bounds}: curves only & L1--L2 \\
\ref{chap07:thm:scaling_law} & Prop. Curves curved in several planes; conditional & text & V & \texttt{Chap07Bounds}: at one point & L1--L2 \\
\ref{chap07:prop:gauss_bonnet_floor} & Prop. Gauss--Bonnet floor for closed surfaces & text & -- & \texttt{Chap07Bounds}: pointwise algebra; conditional on Gauss--Bonnet & L1--L2 \\
\addlinespace[2pt]
\multicolumn{6}{@{}l}{\textbf{Chapter 8. Minimax Curvature and ADM Slicings}} \\[1pt]
\ref{chap08:thm:gauss_codazzi_ricci} & Thm. Equations of Gauss, Codazzi-Mainardi, and Ricci & cited & -- & -- & L1--L2 \\
\ref{chap08:thm:sectional_extrinsic} & Thm. Sectional-Extrinsic Coupling & text & V, F & \texttt{Chap08Pointwise}: pointwise algebra; Gauss equation as hypothesis & L1--L2 \\
\ref{chap08:lem:euclidean_uturn} & Lem. Euclidean U-Turn Bound & text & A & \texttt{Chap08UTurn}: differentiable tangent angle & L1--L2 \\
\ref{chap08:prop:space_form_uturn} & Thm. U-turns in tubes of space forms & text & V, F, A & \texttt{Chap08SpaceForms}: Fermi-coordinate reduction; Clairaut identity as hypothesis & L1--L2 \\
\ref{chap08:thm:adm_constraints} & Thm. ADM Constraint Equations & cited + sketch & -- & -- & L1--L2 \\
\ref{chap08:cor:adm_bounds} & Cor. Pointwise Bounds on Slices of Bounded Extrinsic Curvature & text & V, F & -- & L1--L2 \\
\ref{chap08:prop:raychaudhuri} & Prop. Raychaudhuri Obstruction to Frozen Expansion & text & -- & \texttt{Chap08Pointwise}: along one curve; Raychaudhuri equation as hypothesis & L1--L2 \\
\ref{chap08:prop:rn_horizon} & Prop. Horizon Curvature in Painlev\'e--Gullstrand Slicing & text & V, F & \texttt{Chap08Pointwise}: algebraic core of (2) & L1--L2 \\
\ref{chap08:thm:israel_junction} & Thm. Israel Junction Theorem & cited & -- & -- & L1--L2 \\
\ref{chap08:prop:wormhole_throat} & Prop. Throat Geometry and the Null Energy Bound & text & V, F & \texttt{Chap08Pointwise}: (2)--(3), pointwise at the throat & L1--L2 \\
\ref{chap08:prop:proper_accel} & Prop. Proper Acceleration as Second Fundamental Form & text & V & \texttt{Chap08Pointwise}: linear algebra at one point & L1--L2 \\
\ref{chap08:prop:bh_topology} & Prop. Topology of Multi-Black-Hole Slices & text & -- & -- & L1--L2 \\
\ref{chap08:prop:causal_constraints} & Prop. Causal Constraints on Minimax Worldlines & text & -- & \texttt{Chap08Chronology}: (1)--(2), order-theoretic skeleton & L1--L2 \\
\addlinespace[2pt]
\multicolumn{6}{@{}l}{\textbf{Chapter 9. Global Minimax Curvature}} \\[1pt]
\ref{chap09:prop:holonomy} & Prop. Non-abelian path classification & text & V, A & \texttt{Chap09Covering}: categorical core; faithfulness as hypothesis & L1--L2 \\
\ref{chap09:prop:subloop} & Prop. Confinement of closed sub-loops & text & F & \texttt{Chap09SubloopCorridor}: differentiable tangent angle & L1--L2 \\
\ref{chap09:thm:loop_bounding} & Prop. Finitely many classes under a length bound & text & V & \texttt{Chap09WindingBounds}: (1), one obstacle; homotopy identity as hypothesis & L1--L2 \\
\ref{chap09:prop:winding_total_curvature} & Prop. Winding and total curvature & text & F & \texttt{Chap09Winding}: $C^2$ curves; regularity of the angle lift as hypothesis & L1--L2 \\
\ref{chap09:thm:unfolding} & Prop. Unfolding irreducible immersions & text & V & \texttt{Chap09Covering}: more general than the text & L1--L2 \\
\ref{chap09:prop:corridor_gap} & Prop. No curvature gap in the benchmark & text & F & \texttt{Chap09SubloopCorridor}: lower bound; attaining arcs only & L1--L2 \\
\ref{chap09:prop:teardrop} & Prop. Teardrop U-turn & text & V & \texttt{Chap09Teardrop}: (i) for $a < 2\rho$; (ii) for one turning pattern & L1--L2; F-52 \\
\ref{chap09:prop:mouth_halfturn} & Prop. Mouth width under a half-turn range & text & -- & \texttt{Chap09MouthHalfturn}: differentiable tangent angle in $[0,\pi]$ & L1--L2; F-52 \\
\ref{chap09:prop:balloon_height} & Prop. Height of the balloon & text & -- & \texttt{Chap08UTurn}: differentiable tangent angle & L1--L2; F-52 \\
\ref{chap09:prop:teardrop_lower} & Prop. Lower bound for the balloon radius & text & -- & \texttt{Chap09Teardrop}: conditional, final step only & L1--L2; F-52 \\
\addlinespace[2pt]
\multicolumn{6}{@{}l}{\textbf{Chapter 10. Minimax Trajectories on Statistical Manifolds}} \\[1pt]
\ref{chap10:prop:degenerate} & Prop. Degeneracy & text & A & \texttt{Chap10Generalization}: abstract; geodesic property as hypothesis & L1--L2 \\
\ref{chap10:prop:counterexample} & Prop. Counterexample to curvature-based generalization bounds & text & F, A & \texttt{Chap10Generalization}: rank hypothesis explicit; Gaussian value not formalized & L1--L2; F-59 \\
\addlinespace[2pt]
\multicolumn{6}{@{}l}{\textbf{Chapter 11. Emergent Spacetime and Quantum Geometry}} \\[1pt]
\ref{chap11:prop:relative_entropy_hessian} & Prop. Hessian of relative entropy & text & V & \texttt{Chap11Fisher}: commuting case only & L1--L2 \\
\ref{chap11:thm:first_law_entanglement} & Prop. First law of entanglement & text & V & \texttt{Chap11Fisher}: commuting case only & L1--L2 \\
\ref{chap11:thm:einstein_emergence} & Thm. Faulkner--Guica--Hartman--Myers--Van Raamsdonk & cited + sketch & -- & -- & L1--L2 \\
\ref{chap11:prop:mincut} & Prop. Min-cut bound & text & F & \texttt{Chap11MinCut}: linear-algebra half; factorization as hypothesis & L1--L2 \\
\ref{chap11:prop:mcf_area} & Prop. Area monotonicity under mean curvature flow & text & V & -- & L1--L2 \\
\ref{chap11:thm:spin_net_tensor_equiv} & Prop. Spin networks are contracted tensor networks & cited + proof & -- & -- & L1--L2 \\
\ref{chap11:thm:ashtekar_limit} & Prop. Discrete holonomies and the area spectrum & cited + proof & V, F & -- & L1--L2 \\
\ref{chap11:thm:horizon_complexity} & Prop. Landscape complexity and concentration & cited & -- & -- & L1--L2 \\
\addlinespace[2pt]
\multicolumn{6}{@{}l}{\textbf{Chapter 12. Holographic Entanglement and Quantum Geometry}} \\[1pt]
\ref{chap12:thm:running_spectral_dimension} & Prop. Spectral Dimension of a Homogeneous Symbol & text & -- & \texttt{Chap12Spectral}: heat-kernel diagonal defined by its Fourier formula & L1--L2 \\
\ref{chap12:thm:cartan_metric_emergence} & Prop. Barycentric Hessian and the $A_{m-1}$ Root Metric & text & V & \texttt{Chap12Spectral}: algebraic and calculus core & L1--L2 \\
\ref{chap12:thm:minimax_hamiltonian_regularization} & Thm. Pointwise Constraint Bounds on Slices of Bounded Extrinsic Curvature & text & V & \texttt{Chap12Constraint}, \texttt{Chap12ConstraintMatrix}: pointwise, orthonormal frame; constraint as hypothesis & L1--L2 \\
\ref{chap12:thm:jordan_chronology} & Prop. Timelike Curves Do Not Self-Intersect in Chronological Spacetimes & text & -- & \texttt{Chap08Chronology}: order-theoretic skeleton & L1--L2 \\
\ref{chap12:thm:grand_einstein_emergence} & Thm. Wald Symplectic Equivalence: Linearized Einstein Equations and Second-Order Extension & cited + proof & -- & -- & L1--L2 \\
\ref{chap12:thm:neckpinch_condensation} & Prop. Exponential Decay of Negatively Curved Entries & text & V, A & \texttt{Chap12Spectral}: pointwise ODE, curvature continuous in $t$ & L1--L2; F-59 \\
\addlinespace[2pt]
\end{longtable}
}

\textbf{Chapter 13} contains no theorem-type statements. Its closed form for $d_s(\tau)$ and its numerical bounds ($H_{\mathrm{inf}}$, $|\alpha_t|$, $\ell_\ast$, the rescaled delays) are formalized as arithmetic in \texttt{Chap13Numbers}, with the closed form of the heat-kernel integral and the limits $d_s \to 2, 4$ not covered. Some remarks also have \textsc{Lean} counterparts: the comparison of the SLD and Kubo--Mori metrics in Chapter~11 (\texttt{Chap11Fisher}), the slingshot bound in Chapter~8 (\texttt{Chap09WindingBounds}), the formulation remark of Chapter~10 (\texttt{Chap10Generalization}) and the classification table of Chapter~7 (\texttt{Chap07Bounds}).


\end{document}

